\documentclass[10pt]{article}
\usepackage[margin=0.8cm]{geometry}
\usepackage{amssymb}
\usepackage{pdfpages}
\usepackage{graphbox}
\usepackage{tikz}
\usepackage{multirow} 
\begin{document}

\begin{samepage}
\begin{center}
  {\Large\bf Logarithms to base 10}
\end{center}

\pagestyle{empty}
\thispagestyle{empty}

\centerline{\small
\begin{tabular}{|r|c|ccc|ccc|ccc|ccccccccc|}\hline
  &&&&&&&&&&&  \multicolumn{9}{|c|}{ADD} \\
  {\normalsize x}&
  {\normalsize 0}&{\normalsize 1}&{\normalsize 2}&
  {\normalsize 3}&{\normalsize 4}&{\normalsize 5}&
  {\normalsize 6}&{\normalsize 7}&{\normalsize 8}&
  {\normalsize 9}&{\normalsize 1}&{\normalsize 2}&
  {\normalsize 3}&{\normalsize 4}&{\normalsize 5}&
  {\normalsize 6}&{\normalsize 7}&{\normalsize 8}&
  {\normalsize 9}\\ \hline
\vphantom{$A^{A^A}$}
\input{log_table_values_page1.txt}
\hline
\end{tabular}
}

\noindent{\small For \rule{0mm}{5mm}example, the logarithm of 2.25 is
  0.3522, directly from the table.  To calculate the logarithm
  of~$1.267$, there are two rows for ``1.2''.  The first row is for
  second digits 0-4 and the other row is for second digits 5-9.  The
  second digit is ``6'', so we choose the second row and the entry is
  ``1004''.  To include the third digit, use the column headed ``7''
  in the differences, which gives $1004+24=1028$ to give a result of
  $0.1028$.}

\begin{tikzpicture}[remember picture,overlay]
  \node[anchor=south east,
    xshift=-1cm,
    yshift=1.1cm]
  at (current page.south east){\includegraphics[width=2in]{mono-logo.pdf}};
  \node[anchor=south west,
    xshift=1cm,
    yshift=1.1cm]
  at (current page.south west){{\tt https://github.com/RobinHankin/tables.git}};
\end{tikzpicture}

\end{samepage}
\newpage

\vfill{\ }
\vfill{\ }

\centerline{\small
\begin{tabular}{|r|c|ccc|ccc|ccc|ccccccccc|}\hline
  &&&&&&&&&&&  \multicolumn{9}{|c|}{ADD} \\[1mm]
  {\normalsize x}&
  {\normalsize 0}&{\normalsize 1}&{\normalsize 2}&
  {\normalsize 3}&{\normalsize 4}&{\normalsize 5}&
  {\normalsize 6}&{\normalsize 7}&{\normalsize 8}&
  {\normalsize 9}&{\normalsize 1}&{\normalsize 2}&
  {\normalsize 3}&{\normalsize 4}&{\normalsize 5}&
  {\normalsize 6}&{\normalsize 7}&{\normalsize 8}&
  {\normalsize 9}\\ \hline
  \vphantom{$A^{A^A}$}
  \input{log_table_values_page2.txt}
\hline
\end{tabular}
}

\vfill{\ }
\end{document}

